\documentclass[a4paper,11pt]{article}
\usepackage{exam-style-341}
\examinehead{341 农业综合知识三 · 网络技术分章练习}{第01章 参考答案与解析}
\begin{document}

\unitanswerscover{341 农业综合知识三 · 网络技术分章练习}{第01章参考答案}{计算机网络基本概念与性能指标：逐题答案与解析；计算题给出完整步骤与单位。}

\answerpart{一、单项选择题（10 题，每题 3 分，共 30 分）}

\q A。\begin{analysis}教材 1.1 指出，计算机网络向用户提供的最重要的功能有两个，即连通性和共享（教材第 1 章 1.1，印刷页 1—2）。B、C、D 都不是计算机网络的主要功能。\end{analysis}

\q A。\begin{analysis}互联网基础结构发展分三个阶段：第一阶段从单个网络 ARPANET 向互连网发展；第二阶段建成三级结构（主干网、地区网、校园网）的互联网；第三阶段形成多层次 ISP 结构的互联网（教材 1.2.2，印刷页 5—7）。\end{analysis}

\q A。\begin{analysis}以小写字母 i 开始的 internet（互连网）是通用名词，泛指由多个计算机网络互连而成的网络，协议可以任意选择；以大写字母 I 开始的 Internet（互联网）是专有名词，指采用 TCP/IP 协议族的全球最大特定互连网（教材 1.2.2，印刷页 6—7）。B、C、D 均错。\end{analysis}

\q A。\begin{analysis}按覆盖范围，计算机网络分为广域网 WAN、城域网 MAN、局域网 LAN 和个人区域网 PAN，由大到小依次为广域网、城域网、局域网、个人区域网（教材 1.5，印刷页 19—20）。\end{analysis}

\q A。\begin{analysis}电路交换必须经过“建立连接—通信—释放连接”三个阶段，通信过程中用户始终独占端到端通信资源（教材 1.3.2，印刷页 12—14）。B 是报文交换和分组交换的存储转发过程，C 是 PCM 的三个步骤，D 是分层封装的描述。\end{analysis}

\q D。\begin{analysis}分组交换的主要优点是高效、灵活、迅速、可靠（教材 1.3.2，印刷页 15—16）。D 错误：分组交换在通信前不必占用一条端到端的通信资源，而是在哪段链路上传送才占用那段链路；“独占端到端信道”是电路交换的特点。\end{analysis}

\q A。\begin{analysis}发送时延 $=$ 数据长度（bit）$/$ 发送速率（bit/s），它发生在机器内部的发送器中，与帧长成正比、与发送速率成反比（教材 1.6.1，印刷页 22）。B 是传播时延公式，D 是时延带宽积公式。\end{analysis}

\q C。\begin{analysis}传播时延 $=$ 信道长度 $/$ 电磁波在信道中的传播速率，它发生在机器外部的传输媒体上，与信号发送速率无关（教材 1.6.1，印刷页 22—23）。A、B、D 都与传播时延有关。\end{analysis}

\q A。\begin{analysis}时延带宽积 $=$ 传播时延 $\times$ 带宽，表示链路上可容纳的比特数，又称“以比特为单位的链路长度”（教材 1.6.1，印刷页 24—25）。\end{analysis}

\q B。\begin{analysis}RTT 是从发送方发送数据开始、到发送方收到接收方确认所经历的时间；互联网中的 RTT 还包括各中间节点的处理时延、排队时延以及转发数据时的发送时延（教材 1.6.1，印刷页 25）。A 只取了一部分，C、D 均错。\end{analysis}

\answerpart{二、多项选择题（3 题，每题 4 分，共 12 分）}

\q ABC。\begin{analysis}教材 1.6.1 介绍常用的 7 个性能指标：速率、带宽、吞吐量、时延、时延带宽积、往返时间 RTT 和利用率。D 显示器分辨率属于显示设备指标，不是网络性能指标。\end{analysis}

\q ABC。\begin{analysis}A 正确：电路交换通信前必须先建立连接；B 正确：报文交换以整个报文为传输单位；C 正确：分组交换的每个分组携带首部，可独立选择传输路径（教材 1.3.2）。D 错误：电路交换线路利用率往往很低，不适合传送突发性的计算机数据。\end{analysis}

\q ABC。\begin{analysis}A、B 正确：带宽有频域（Hz）和时域（bit/s）两种意义；C 正确：吞吐量是单位时间内实际通过的数据量；D 错误：吞吐量受带宽、网络负载和拥塞等影响，通常小于带宽，并不总是等于带宽（教材 1.6.1）。\end{analysis}

\answerpart{三、判断题（6 题，每题 2 分，共 12 分）}

\q ×。改正：计算机网络中的资源包括硬件资源、软件资源和数据资源。\begin{analysis}教材 1.1 指出，计算机网络中的资源包括硬件、软件和数据，资源共享是计算机网络的重要功能之一。\end{analysis}

\q √。\begin{analysis}小写 internet 是通用名词，泛指由多个计算机网络互连而成的网络，网络之间的通信协议可以任意选择（教材 1.2.2，印刷页 6—7）。\end{analysis}

\q √。\begin{analysis}电路交换在通信过程中独占端到端的通信资源，时延小；但计算机数据具有突发性，线路空闲时资源被浪费，因此线路传输效率往往很低（教材 1.3.2，印刷页 12—14）。\end{analysis}

\q √。\begin{analysis}分组交换采用存储转发方式，分组在各路由器需要先排队、查表再转发，因此会带来排队时延；通信量过大时排队时延会显著增大（教材 1.3.2、1.6.1）。\end{analysis}

\q ×。改正：发送时延与数据长度和发送速率有关，与信道长度无关；传播时延与信道长度和传播速率有关，与发送速率无关。\begin{analysis}教材 1.6.1（印刷页 22—23）把两种时延的区别作为重点：发送时延发生在机器内部发送器，传播时延发生在传输媒体上。\end{analysis}

\q ×。改正：信道利用率并非越高越好，利用率过高时时延会迅速增加，一般控制信道利用率不超过 50\%。\begin{analysis}教材 1.6.1（印刷页 25—26）根据排队论指出，当网络利用率接近最大值时，网络时延趋于无穷大，因此 ISP 通常控制信道利用率不超过 50\%（式 1-5）。\end{analysis}

\answerpart{四、计算题（1 题，共 6 分）}

\q 参考答案：
（1）发送时延 $=$ 数据长度 $/$ 发送速率 $=5000\ \text{bit}\div(10\times10^{6}\ \text{bit/s})=5\times10^{-4}\ \text{s}=0.5\ \text{ms}$（1.5 分）。
（2）传播时延 $=$ 信道长度 $/$ 传播速率 $=2000\ \text{km}\div(2\times10^{8}\ \text{m/s})=2\times10^{6}\ \text{m}\div(2\times10^{8}\ \text{m/s})=10^{-2}\ \text{s}=10\ \text{ms}$（1.5 分）。
（3）RTT $\approx 2\times$ 单向传播时延 $=2\times10\ \text{ms}=20\ \text{ms}$；确认帧长度忽略不计，若计入处理和排队时延，RTT 还会更大（1.5 分）。
（4）时延带宽积 $=$ 传播时延 $\times$ 带宽 $=10\times10^{-3}\ \text{s}\times10\times10^{6}\ \text{bit/s}=1\times10^{5}\ \text{bit}=100\ \text{kbit}=12.5\ \text{KB}$（1.5 分）。
\begin{analysis}本题只考查发送时延、传播时延、RTT 和时延带宽积四个概念。注意两点：一是发送时延与链路长度无关，传播时延与发送速率无关；二是 RTT 严格说还应包括确认帧的发送时延、处理和排队时延，本题已说明忽略。\end{analysis}

\answerpart{五、简答题（2 题，每题 10 分，共 20 分）}

\q 参考答案要点：
（1）定义（3 分）：计算机网络是把地理位置不同、具有独立功能的多台计算机及其外部设备，通过通信线路连接起来，在网络协议和网络软件的管理协调下，实现资源共享和信息传递的计算机系统。
（2）两大功能（3 分）：连通性——使上网用户之间可以交换信息，好像这些用户的计算终端直接相连；共享——包括硬件共享、软件共享和数据（信息）共享。
（3）三种交换方式（4 分）：电路交换在通信前建立连接、通信中独占端到端资源、通信后释放连接，时延小但线路利用率低；报文交换以整个报文为单位存储转发，无需建立连接，但时延大、对节点存储要求高；分组交换把报文分组，以分组为单位存储转发、独立选路，具有高效、灵活、迅速、可靠的优点，但有首部开销、可能排队和失序。（教材 1.3.2，印刷页 12—16）

\q 参考答案要点（每个指标 1 分，时延四部分共 3 分）：
（1）速率：数据的传送速率，单位 bit/s（教材 1.6.1）。
（2）带宽：频域指频带宽度（Hz），时域指最高数据率（bit/s）。
（3）吞吐量：单位时间内通过某网络或信道的实际数据量，受带宽和负载限制。
（4）时延：数据从链路一端传送到另一端所需的时间，由四部分组成——发送时延 $=$ 数据长度 $/$ 发送速率；传播时延 $=$ 信道长度 $/$ 传播速率；处理时延（分析首部、查表、差错检验等）；排队时延（在路由器输入／输出队列等待），总时延为四者之和。
（5）时延带宽积 $=$ 传播时延 $\times$ 带宽，表示链路上可容纳的比特数。
（6）往返时间 RTT：从发送数据开始到收到确认所经历的时间。
（7）利用率：信道利用率与网络利用率，利用率过高时时延迅速增大。
其中，发送时延与帧长、发送速率有关；传播时延与距离、传播速率有关；处理时延取决于设备处理能力；排队时延取决于网络当前通信量。（教材 1.6.1，印刷页 21—26）

\answerpart{六、综合题（1 题，共 20 分）}

\q 参考答案要点：
（1）比特数与发送时延（4 分）：$200\ \text{MB}=200\times10^{6}\ \text{B}=200\times10^{6}\times8\ \text{bit}=1.6\times10^{9}\ \text{bit}$；发送时延 $=1.6\times10^{9}\ \text{bit}\div(100\times10^{6}\ \text{bit/s})=16\ \text{s}$。
（2）判断（4 分）：不严谨。16 s 只是链路带宽为 100 Mb/s 时的发送时延，实际传输时间还取决于传播时延、处理和排队时延、协议开销、服务器与云端的处理速度以及链路竞争；实际吞吐量通常小于带宽，因此文件传完的时间一般大于 16 s。
（3）传播时延与时延带宽积（6 分）：传播时延 $=1000\ \text{km}\div(2\times10^{8}\ \text{m/s})=10^{6}\ \text{m}\div(2\times10^{8}\ \text{m/s})=5\times10^{-3}\ \text{s}=5\ \text{ms}$；时延带宽积 $=5\ \text{ms}\times100\ \text{Mb/s}=5\times10^{-3}\times100\times10^{6}\ \text{bit}=5\times10^{5}\ \text{bit}=500\ \text{kbit}=62.5\ \text{KB}$。物理意义：表示该链路像一条管道，最多可容纳约 500 kbit（62.5 KB）正在传输而尚未到达接收端的比特。
（4）改进建议（6 分，答出两条即可）：（1） 提高出口链路带宽，如升级链路或采用链路聚合；（2） 优化排队与调度、为视频业务配置 QoS 优先级，减少排队时延；（3） 采用内容分发网络（CDN）或就近上传，缩短传输距离、降低传播时延；（4） 对视频进行压缩或分块并行传输，降低数据量、提高吞吐量；（5） 选择网络空闲时段传输，减少竞争。

\end{document}
